\chapter{Conclusão} \label{sec:conclusion}

Este trabalho avaliou ganhos de proficiência associados à adesão ao Programa Ensino Integral nas escolas estaduais paulistas no SARESP entre 2011 e 2018. A pergunta central não foi apenas se as escolas PEI melhoraram após a conversão, mas se esse ganho permanece quando se leva a sério o desenho empírico do problema: adoção escalonada, efeitos heterogêneos no tempo, diferenças observáveis entre escolas e mudança composicional do alunado. As estimativas principais indicam efeitos médios pós-tratamento entre aproximadamente 0,50 e 0,68 desvio-padrão, com efeitos maiores em Matemática e no Ensino Médio. A dinâmica crescente dos \textit{event studies} poderia refletir seleção dinâmica de alunos, mudanças de participação na prova ou recuperação de escolas inicialmente abaixo da média; também é compatível com maturação pedagógica e organizacional do programa. Os dados disponíveis não permitem separar completamente esses canais.

A contribuição substantiva do trabalho é mostrar que esses resultados não dependem apenas de uma comparação simples entre escolas tratadas e controles. Os ganhos permanecem quando se abandona o TWFE como estimador principal, quando se usa \textcite{callawaysantanna} para lidar com a adoção escalonada, quando se adicionam covariáveis por meio de \textcite{santannazhao} e quando se implementa uma aproximação agregada inspirada em \textcite{santannaxu}. Essa última etapa é importante porque a literatura sobre o PEI mostra que a conversão das escolas altera a composição dos alunos. Portanto, uma avaliação do programa que ignore esse canal corre o risco de confundir ganho de aprendizagem com seleção discente. Dentro das dimensões observáveis disponíveis neste trabalho, a correção composicional não elimina os efeitos estimados, mas o diagnóstico tipo Hausman tem baixa potência e não constitui evidência contra efeitos composicionais de até aproximadamente 0,20--0,25 desvio-padrão.

A contribuição metodológica é explicitar a distância entre o estimador ideal e o estimador possível. O estimador de \textcite{santannaxu} foi desenhado para cortes transversais repetidos em que a mudança composicional pode ser tratada no nível individual. Para aplicá-lo de forma completa, seria necessário observar, no mesmo registro, a matrícula individual, as covariáveis do aluno e sua nota no SARESP. Essa estrutura não está disponível publicamente para o período analisado. Por isso, este trabalho construiu uma adaptação agregada: para cada coorte e tempo relativo, foram estimadas células 2x2 com logit multinomial de quatro categorias, pesos de composição e regressões de resultado, sempre usando covariáveis fixadas no baseline da coorte. A versão final não usa fallback binário para S\&Z; quando não há suporte comum suficiente, a célula é excluída. Isso torna a estimativa mais restrita, mas também mais honesta.

Essa limitação define o alcance da interpretação. Os resultados são evidência de ganhos agregados de proficiência nas escolas PEI; não são uma decomposição perfeita entre aprendizado individual, seleção de alunos e mudanças de participação na prova. Ainda assim, o exercício é informativo. Se a seleção observável fosse a principal explicação dos ganhos, seria esperado que as estimativas caíssem de forma importante ao passar do CS21 para o DR-DiD e para a aproximação S\&X. Não é isso que ocorre. Além disso, os diagnósticos de regressão à média e de spillovers regionais preservam magnitudes próximas às principais. Por outro lado, a sensibilidade sem imputação mostra que a magnitude depende da cobertura das covariáveis de matrícula, especialmente porque a amostra com dados completos é muito menor e, no Ensino Médio, insuficiente para sustentar a mesma estimação. O novo diagnóstico de participação no SARESP também recomenda cautela adicional no Ensino Médio, onde a razão entre participantes e matrículas cai após a conversão.

As principais pontas abertas decorrem dessa mesma restrição de dados. A primeira extensão natural é obter acesso a bases administrativas que permitam vincular matrícula individual, participação na prova e nota individual, aproximando de forma mais fiel o estimador de Sant'Anna e Xu e separando melhor aprendizado, seleção e comparecimento ao exame. A segunda é incorporar medidas geográficas mais diretas de exposição a escolas PEI vizinhas, usando distância ou tempo de deslocamento, para avançar além da proxy por diretoria de ensino e distrito. A terceira é expandir a janela para anos posteriores a 2018, tratando separadamente a pandemia, as mudanças de avaliação e o crescimento acelerado do PEI. Em suma, este trabalho mostra que o PEI está associado a ganhos persistentes de proficiência em um painel agregado de escolas. A conclusão não encerra a discussão sobre mecanismos, mas organiza melhor seus termos: há ganho agregado positivo, há mudança composicional, há indícios de alteração na participação no Ensino Médio, e a parte observável da composição não parece suficiente para derrubar os resultados sob as especificações estimadas. Ao mesmo tempo, as magnitudes obtidas são superiores às estimadas por \textcite{scorzafave2023} para o mesmo programa, o que reforça que a leitura mais prudente é tratar os resultados como evidência agregada condicionada ao desenho disponível, e não como uma medida definitiva do efeito individual líquido do PEI.
